\documentclass[a4paper, 12pt]{article}

\usepackage[french]{babel}
\usepackage[T1]{fontenc}
\usepackage[utf8]{inputenc}
\usepackage{amsmath, amssymb}
\usepackage{geometry}
\usepackage{array}
\usepackage{enumitem}
\usepackage{tikz}
\usepackage{tkz-tab}
\usepackage{pgfplots}
\pgfplotsset{compat=1.18}

\geometry{top=2cm, bottom=2cm, left=2cm, right=2cm}

\begin{document}

\begin{center}
    \Large\textbf{Corrigé du Bac Blanc - Spécialité Mathématiques}
\end{center}

\vspace{0.5cm}

\section*{Partie 1 : Automatismes (6 points)}

\begin{enumerate}[leftmargin=*, itemsep=0.3cm]
    \item \textbf{Question 1 :} $\frac{1}{3}-\frac{2-x}{2} = \frac{2-6+3x}{6} = \frac{-4+3x}{6}$. \textbf{Réponse d}.
    \item \textbf{Question 2 :} $B = \frac{1/2}{1/2} + \frac{1/6}{1/3} = 1 + \frac{1}{2} = \frac{3}{2}$. \textbf{Réponse b}.
    \item \textbf{Question 3 :} $x = \frac{2y}{Cy-3}$ (en isolant $x$). \textbf{Réponse b}.
    \item \textbf{Question 4 :} $P = \frac{1}{4} \text{ (majeurs)} \times \frac{2}{3} \text{ (18-25 ans)} = \frac{1}{6}$. \textbf{Réponse a}.
    \item \textbf{Question 5 :} $36 \text{ m}^3\text{/h} = \frac{36000 \text{ L}}{3600 \text{ s}} = 10 \text{ L/s}$. \textbf{Réponse b}.
    \item \textbf{Question 6 :} $f(x) \le 5 \iff x^2 \ge 5$, soit $x \in ]-\infty ; -\sqrt{5}] \cup [\sqrt{5} ; +\infty[$. \textbf{Réponse b}.
    \item \textbf{Question 7 :} $\frac{10+14+2X}{5} = 12 \implies 24+2X = 60 \implies X = 18$. \textbf{Réponse b}.
    \item \textbf{Question 8 :} $CM = 0,7 \times 0,8 = 0,56$, d'où une baisse de $44\%$. \textbf{Réponse c}.
\end{enumerate}

\section*{Partie 2 : Exercices (14 points)}

\subsection*{Exercice 1 : Étude de fonctions et suites}

\textbf{Partie A : Modèle continu}
\begin{enumerate}
    \item $g(x)=f(x) \iff \frac{3+x}{2x} = x-2 \iff \frac{-2x^2+5x+3}{2x} = 0$.
    \item $-2x^2+5x+3=0 \implies \Delta = 49 \implies x_1 = 3$ et $x_2 = -0,5$.
    Points d'intersection : $I_1(3 ; 1)$ et $I_2(-0,5 ; -2,5)$.
    \item \textbf{Tableau de signes de $f(x) = \frac{-2x^2+5x+3}{2x}$ :}
    \begin{center}
    \begin{tikzpicture}
        \tkzTabInit[lgt=4, espcl=1.5]{$x$ / 1, $-2x^2+5x+3$ / 1, $2x$ / 1, $f(x)$ / 1.5}
        {$-\infty$, $-0.5$, $0$, $3$, $+\infty$}
        \tkzTabLine{ , - , 0 , + , t , + , 0 , - }
        \tkzTabLine{ , - , t , - , d , + , t , + }
        \tkzTabLine{ , + , 0 , - , d , + , 0 , - }
    \end{tikzpicture}
    \end{center}
\end{enumerate}

\textbf{Partie B : Modèle discret}
\begin{enumerate}
    \item $u_2 = 0$ et $u_3 = 1$. $v_2 = 1,25$ et $v_4 = 0,875$.
    \item $u_{n+1}-u_n = 1 > 0$ (croissante). $v_{n+1}-v_n = \frac{-3}{2n(n+1)} < 0$ (décroissante).
    \item $w_{n+1}-w_n = (v_{n+1}-v_n) - (u_{n+1}-u_n) < 0$. La suite $(w_n)$ est décroissante.
\end{enumerate}

\subsection*{Exercice 2 : Mouvement du piston}
$v_{max} = 0,65 \text{ m/s}$ (à $t = \frac{3\pi}{26} \text{ s}$) ; $v_{min} = -0,65 \text{ m/s}$ (à $t = \frac{\pi}{26} \text{ s}$).
Vitesse nulle pour $t \in \{0 ; \frac{\pi}{13} ; \frac{2\pi}{13}\}$.

\subsection*{Exercice 3 : Géométrie et produit scalaire}
\begin{enumerate}
    \item En utilisant les coordonnées $E(0,0), K(1,2), M(3,1)$ : $\vec{EK} \cdot \vec{EM} = 1 \times 3 + 2 \times 1 = 5$.
    \item $EL = \frac{\vec{EK} \cdot \vec{EM}}{EM} = \frac{5}{\sqrt{10}} = \frac{\sqrt{10}}{2}$.
    \item $\cos(\widehat{KEM}) = \frac{5}{\sqrt{5}\sqrt{10}} = \frac{1}{\sqrt{2}} \implies \widehat{KEM} = 45^\circ$.
    \item $\vec{OA} \cdot \vec{OB} = x^2+2x-6$. Minimum atteint pour $x = -1$.
    \item Pour $x=-1$, $\cos(\widehat{BOA}) = \frac{-7}{\sqrt{65}} \implies \widehat{BOA} \approx 150^\circ$.
\end{enumerate}

\end{document}